% ---
% id: receiver/detector-imperfections/rem:normalization-scope
% kind: rem
% title: "What the normalization factors out and what it leaves in place"
% summary: >-
%   Only the per-port efficiencies are divided: declared mode-dependent loss
%   stays, undeclared efficiency mismatch is outside the assumptions either
%   way, and a probe-fed receiver is not covered. The full normalization can
%   remove the loss the noise floor needs, which a larger factor keeps.
% prominence: standard
% status: accepted
% uses: [prop:efficiency-normalization, def:normalized-station, thm:efficiency, prop:noise-floor]
% mentions: [receiver/ass:modes, receiver/ass:port, rem:powervsclick, req:efficiency-normalization]
% assumes: []
% refs: []
% provenance: newly-derived
% lean:
% notes: >-
%   The engine's handling at the full normalization, which is the reason the
%   default factor of req:efficiency-normalization keeps a margin below
%   99/100, is a property of the code that no claim establishes; it is
%   stated in the receiver's design.md (structural decision 14), not in this
%   remark, and is revised there once the engine handles a singular block
%   at the full normalization. On a BBM92 point at efficiency 0.8 a run lost
%   part of the key at t = 0.80001 (about 1e-5 relative above eta_max), all
%   of it at t = 0.800001, and none at t = 0.8001 (about 1e-4).
% ---

\begin{remark}\label{rem:normalization-scope}
\emph{Only the terminal scalars are divided.} The family $\{A^{\circ}_k\}$ of
Definition~\ref{def:normalized-station} is kept: every loss declared in the
element graph, a polarizer's extinction, and an attenuation that varies
within a port (Theorem~\ref{thm:efficiency}\,(ii)) stay in the normalized
station unchanged. A detector whose efficiency depends on a spatial, spectral
or temporal mode is declared as such a mode-dependent loss in front of its
port, and the normalization leaves it where it is.

\emph{Efficiency mismatch is a separate problem.} The identity of
Proposition~\ref{prop:efficiency-normalization}\,(i) holds of the declared
instrument. That it describes the physical receiver, clause (iii), rests on
Assumptions~\ref{ass:modes} and~\ref{ass:port}: an efficiency that differs
between degrees of freedom the adversary can address but the declaration
does not carry violates one of them, with or without the normalization, which
neither creates such a mismatch nor removes one.

\emph{A probe-fed receiver is not covered.} With declared source modes the
probe light passes the same detectors, so a uniform loss moved to the front
would act on the probe modes too and replace each probe state by its image
under the loss; no claim here states that form.

\emph{The choice of the factor.} The identity holds at every
$t\in[\eta_{\max},1]$. At the full normalization $t=\eta_{\max}$ with lossless
optics the normalized station can have $\lmax(A'_\Sigma)=1$, and then some
element of its dark-composed family is singular on every block with a photon,
dark counts or not (Proposition~\ref{prop:noise-floor}\,(iv)). At
$t>\eta_{\max}$ the largest efficiency of the station is $\eta_{\max}/t<1$,
so $\lmax(A'_\Sigma)\le\eta_{\max}/t$ and, with a dark count in every port,
every such element is positive definite with the floor of
Proposition~\ref{prop:noise-floor}\,(iii) at $\eta_{\max}/t$. Where
$\eta_{\max}<99/100$, the default factor
$t=\eta_{\max}/(99/100)$ of Requirement~\ref{req:efficiency-normalization}
is of this kind, with $\lmax(A'_\Sigma)\le99/100$; where
$\eta_{\max}\ge99/100$ the default is $t=1$, the declared receiver, whose
largest efficiency can itself be one. Which factor in $[\eta_{\max},1]$ a run
is better served by is a matter of the engine's evaluation, not of the
identity.

\emph{Only declared efficiencies are factored.} The proof of
Proposition~\ref{prop:efficiency-normalization}\,(i) uses only $A_k=t\,A'_k$
with $\{A'_k\}$ admissible, which holds for every
$t\in[\lmax(A_\Sigma),1]$ with $t>0$ and $A'_k=A_k/t$; such a factor can be smaller
than $\eta_{\max}$, but it is an eigenvalue rather than a declared number.
A Gram family declared from click data, with every efficiency declared as one
(Remark~\ref{rem:powervsclick}), has $\eta_{\max}=1$, and nothing is factored
from it.
\end{remark}
